% SHARD-ID: AQM-48-F3-FULL-AFFINE-LINE-NET-HILBERT-ACTION
% SHARD-TITLE: The Full \(\mathbb F_3\) Affine-Line Net and Hilbert Action
% SHARD-SUMMARY: Separates the rational three-site truncation from the full countable closed-point net on \(\mathbf A^1_{\mathbb F_3}\).
% SHARD-SUMMARY: Builds the zero-reference incomplete tensor Hilbert representation and proves the affine symmetry implementation.
% SHARD-SUMMARY: Identifies the generic point as fixed but not part of the finite-residue Weyl net.
% SHARD-KEYWORDS: F3, affine line, full net, incomplete tensor product, generic point, Clifford

\section{The Full \texorpdfstring{\(\mathbb F_3\)}{F3} Affine-Line Net and Hilbert Action}
\label{sec:f3-full-affine-line-net-hilbert-action}

\subsection{Status and Source Anchors}
This shard corrects the scope of AQM-47.  The three-rational-point calculation
there is a finite closed-support truncation.  It is not the full local net on
\(\mathbf A^1_{\mathbb F_3}\).  The full closed-point net uses every monic
irreducible polynomial in \(\mathbb F_3[x]\).  AQM-28 supplies the
closed-point quasi-local convention, AQM-29 supplies the point and open-set
calculation for \(\Spec k[t]\), AQM-45 supplies the reference-vector incomplete
tensor representation pattern, and AQM-46 supplies the affine symmetry action.

\subsection{Lamport Dependency Skeleton}
\[
\begin{array}{rcl}
D1&:&k=\mathbb F_3,\quad R=k[x],\quad X=\Spec R.\\
D2&:&\eta=(0),\quad \mathcal P=\{\pi\in k[x]:\pi
      \text{ monic irreducible}\}.\\
L1&:&X=\{\eta\}\cup\{x_\pi=(\pi):\pi\in\mathcal P\},\quad
      \mathcal P\text{ is countably infinite}.\\
D3&:&K_\pi=R/(\pi),\quad \mathcal H_\pi=\ell^2(K_\pi),\quad
      \mathcal A_\pi=\operatorname{End}(\mathcal H_\pi).\\
D4&:&E(U)=\bigoplus_{x_\pi\in U}K_\pi^2,\quad
      \mathcal A(U)=\varinjlim_{S\Subset C(U)}
      \bigotimes_{x_\pi\in S}\mathcal A_\pi.\\
L2&:&\text{Every nonempty open }U\text{ has cofinite }
      C(U)\subset\mathcal P.\\
T1&:&\mathcal A(X)\text{ is an algebraic infinite tensor product over }
      \mathcal P.\\
D5&:&\Gamma_{\rm fin}=\bigoplus_{\pi\in\mathcal P}K_\pi,\quad
      \mathcal H^0=\ell^2(\Gamma_{\rm fin}).\\
L3&:&\mathcal H^0\text{ is countably infinite-dimensional and represents }
      \mathcal A(X).\\
D6&:&G=\{g_{a,b}:x\mapsto ax+b,\ a\in\mathbb F_3^\times,\ b\in\mathbb F_3\}.\\
D7&:&(V_gs)_{\pi^g}=\tau_{g,\pi}^{-1}(s_\pi),\quad
      U_g|s\rangle=|V_gs\rangle.\\
T2&:&U_g\mathcal A(U)U_g^*=\mathcal A(gU),\quad
      U_gW_U(e)U_g^*=W_{gU}(S_ge).\\
L4&:&g(\eta)=\eta,\quad \eta\in U\text{ for every nonempty open }U.\\
T3&:&\eta\text{ is not a finite-residue qudit site in this net.}
\end{array}
\]

\subsection{The Countable Closed-Point Set D1--L1}
Let
\[
  R=\mathbb F_3[x],\qquad X=\Spec R.
\]
The zero ideal is prime because \(R\) is a domain; write
\[
  \eta=(0).
\]
For each monic irreducible \(\pi\in R\), write
\[
  x_\pi=(\pi).
\]

\paragraph{Lemma \(L1\).}
The points of \(X\) are exactly
\[
  \{\eta\}\cup\{x_\pi:\pi\in\mathcal P\}.
\]
The set \(\mathcal P\) is countably infinite.

\emph{Proof.}
The point classification is AQM-29 with \(k=\mathbb F_3\).  Directly:
\(\mathbb F_3[x]\) is Euclidean, so every nonzero prime ideal has a monic
generator.  The generator must be irreducible, and every monic irreducible
generates a maximal, hence prime, ideal.

For countability, \(R\) is the union over degrees \(n\ge0\) of the finite sets
of polynomials of degree at most \(n\).  Hence \(R\), and therefore
\(\mathcal P\), is countable.  For infinitude, suppose
\(\pi_1,\ldots,\pi_m\) were all monic irreducibles.  Then
\[
  F=\pi_1\cdots\pi_m+1
\]
is nonconstant.  Choose a monic nonconstant divisor \(\rho\) of \(F\) with
least degree; then \(\rho\) is irreducible.  No \(\pi_i\) divides \(F\),
because \(\pi_i\) divides \(F-1\).  Thus \(\rho\) is a new monic irreducible, a
contradiction.

\subsection{The Full Closed-Point Net D3--T1}
At the closed point \(x_\pi\), put
\[
  K_\pi=R/(\pi).
\]
If \(\deg\pi=d\), then \(K_\pi\) has \(3^d\) elements.  Define
\[
  \mathcal H_\pi=\ell^2(K_\pi),\qquad
  \mathcal A_\pi=\operatorname{End}_{\mathbb C}(\mathcal H_\pi)
  \cong M_{3^d}(\mathbb C).
\]
For a finite closed support \(S\Subset\mathcal P\), set
\[
  \mathcal A(S)=\bigotimes_{\pi\in S}\mathcal A_\pi.
\]

For an open \(U\subset X\), define
\[
  C(U)=\{\pi\in\mathcal P:x_\pi\in U\}.
\]
The label group and local algebra are
\[
  E(U)=\bigoplus_{\pi\in C(U)}K_\pi^2,
  \qquad
  \mathcal A(U)=
  \varinjlim_{S\Subset C(U)}\mathcal A(S),
\]
where the transition maps are \(a\mapsto a\otimes\mathbf 1\).

\paragraph{Lemma \(L2\).}
If \(U\ne\emptyset\), then \(C(U)\) is cofinite in \(\mathcal P\).  If
\(U=\emptyset\), then \(C(U)=\emptyset\).

\emph{Proof.}
Every closed subset of \(\Spec R\) has the form \(V(I)\).  Since \(R\) is a
principal ideal domain, \(I=(h)\).  Thus every open is \(D(h)\).  If \(h=0\),
then \(D(h)=\emptyset\).  If \(h\ne0\), then
\[
  x_\pi\notin D(h)\quad\Longleftrightarrow\quad h\in(\pi)
  \quad\Longleftrightarrow\quad \pi\mid h.
\]
A nonzero polynomial has only finitely many monic irreducible divisors.  Hence
the closed points omitted by \(D(h)\) are finite.

\paragraph{Theorem \(T1\).}
The global closed-point algebra is
\[
  \mathcal A(X)=
  \varinjlim_{S\Subset\mathcal P}
  \bigotimes_{\pi\in S}M_{3^{\deg\pi}}(\mathbb C).
\]
It is an algebraic infinite tensor product over all closed points.  The
rational three-site algebra of AQM-47 is only the finite subalgebra for
\[
  S_{\rm rat}=\{x,x-1,x-2\}\subset\mathcal P.
\]

\emph{Proof.}
For \(X=D(1)\), \(C(X)=\mathcal P\), so the displayed formula is Definition
\(D4\).  Since \(\mathcal P\) is infinite by Lemma \(L1\), this is not a finite
matrix algebra.  The rational support is a finite subset of \(\mathcal P\), so
\(\mathcal A(S_{\rm rat})\) maps into the filtered colimit.  It does not equal
the colimit, because any irreducible \(\pi\notin S_{\rm rat}\) contributes the
nontrivial one-site algebra \(\mathcal A_\pi\).

\subsection{The Infinite Hilbert Representation D5--L3}
The algebraic net above is defined before choosing any Hilbert-space
completion.  To get the standard reference representation, choose the reference
vector \(|0\rangle_\pi\in\mathcal H_\pi\) at every closed point.  Define
\[
  \Gamma_{\rm fin}
  =
  \bigoplus_{\pi\in\mathcal P}K_\pi
  =
  \{s=(s_\pi)_\pi:s_\pi=0\text{ for all but finitely many }\pi\}.
\]
Set
\[
  \mathcal H^0=\ell^2(\Gamma_{\rm fin}),\qquad
  |s\rangle=\bigotimes_{\pi\in\mathcal P}|s_\pi\rangle,
\]
with all but finitely many tensor factors equal to \(|0\rangle_\pi\).

\paragraph{Lemma \(L3\).}
\(\mathcal H^0\) is countably infinite-dimensional, and it represents
\(\mathcal A(X)\).

\emph{Proof.}
The set of finite subsets of the countable set \(\mathcal P\) is countable.
For a fixed finite support \(S\), the set \(\prod_{\pi\in S}K_\pi\) is finite.
Thus \(\Gamma_{\rm fin}\), a countable union of finite sets, is countable.
It is infinite because the configurations supported at one closed point with
value \(1\) are distinct as \(\pi\) varies.  Hence
\(\ell^2(\Gamma_{\rm fin})\) has a countably infinite orthonormal basis.

For finite \(S\), \(\mathcal A(S)\) acts on the coordinates in \(S\) and fixes
all other coordinates.  If \(S\subset T\), this action agrees with the action
of \(a\otimes\mathbf 1_{\mathcal H(T\setminus S)}\).  Therefore the compatible
finite-support actions define a representation of the filtered colimit
\(\mathcal A(X)\).

\subsection{The Full Affine Symmetry Implementation D6--T2}
Let
\[
  G=\operatorname{Aut}_{\mathbb F_3}(X)
   =\{g_{a,b}:x\mapsto ax+b,\ a\in\mathbb F_3^\times,\ b\in\mathbb F_3\}.
\]
For \(\pi\in\mathcal P\), AQM-46 gives
\[
  g(x_\pi)=x_{\pi^g},\qquad
  \pi^g(X)=a^{\deg\pi}\pi((X-b)/a),
\]
and a residue-field isomorphism
\[
  \tau_{g,\pi}:K_{\pi^g}\to K_\pi,\qquad
  \tau_{g,\pi}([X])=a[X]+b.
\]
Define a map on finite-support configurations by
\[
  (V_gs)_{\pi^g}=\tau_{g,\pi}^{-1}(s_\pi).
\]
Because \(\tau_{g,\pi}^{-1}(0)=0\), \(V_g\) sends finite-support configurations
to finite-support configurations.  It is bijective, with inverse
\(V_{g^{-1}}\).  Hence
\[
  U_g|s\rangle=|V_gs\rangle
\]
defines a unitary on \(\mathcal H^0\).

\paragraph{Theorem \(T2\).}
For every open \(U\subset X\) and every \(e\in E(U)\),
\[
  U_gW_U(e)U_g^*=W_{gU}(S_ge).
\]
Consequently
\[
  U_g\mathcal A(U)U_g^*=\mathcal A(gU),
\]
and \(g\mapsto U_g\) is a unitary implementation of the full closed-point
net action in the zero-reference representation.

\emph{Proof.}
At one closed point, the basis relabelling
\(|s\rangle\mapsto|\tau_{g,\pi}^{-1}s\rangle\) carries translation by \(q\) to
translation by \(\tau_{g,\pi}^{-1}q\).  It carries the phase labelled by \(p\)
to the phase labelled by \(\tau_{g,\pi}^{-1}p\), because AQM-46 proves
absolute-trace preservation for \(\tau_{g,\pi}\).  Thus the formula holds on a
one-site Weyl operator.  Tensoring over the finite support of \(e\) proves the
displayed Weyl formula.  The Weyl operators span each finite-support algebra,
so the algebra statement follows.  The maps \(V_g\) compose because the residue
maps \(\tau_{g,\pi}\) compose, again by AQM-46; hence the \(U_g\) compose on the
basis of \(\mathcal H^0\).

\subsection{The Generic Point L4--T3}
\paragraph{Lemma \(L4\).}
The generic point \(\eta=(0)\) is fixed by every \(g\in G\), and every nonempty
open \(U\subset X\) contains \(\eta\).

\emph{Proof.}
The coordinate pullback of \(g\) is an automorphism of the domain
\(\mathbb F_3[x]\), so the inverse image of the zero ideal is the zero ideal.
Thus \(g(\eta)=\eta\).  By Lemma \(L2\), every nonempty open is \(D(h)\) with
\(h\ne0\).  Since \(h\notin(0)\), the point \(\eta\) lies in \(D(h)\).

\paragraph{Theorem \(T3\).}
The generic point is not a tensor factor of the finite-residue Weyl net above.

\emph{Proof.}
Its residue field is
\[
  \kappa(\eta)=\mathbb F_3(x),
\]
the fraction field of \(\mathbb F_3[x]\).  This field is infinite.  Convention
(y) attaches finite Weyl sites only to closed points, whose residue fields are
finite.  Therefore \(\eta\) affects the topology of opens, but it contributes
no factor \(K_\eta^2\), no finite qudit \(\ell^2(K_\eta)\), and no coordinate
of \(\Gamma_{\rm fin}\) in the closed-point net.

If one adds the rational-function sector of convention (z), an additional
choice of additive character is required.  The closed-point rule
\((q,p)\mapsto(\tau^{-1}q,\tau^{-1}p)\) must not be copied blindly.  For the
default functional \(\lambda\) that extracts the coefficient of \(x^{-1}\) at
infinity, let \(\rho:x\mapsto2x\).  Then \(\tau_\rho^{-1}(x^{-1})=2x^{-1}\), so
\[
  \lambda(1\cdot x^{-1})=1,\qquad
  \lambda(\tau_\rho^{-1}(1)\tau_\rho^{-1}(x^{-1}))=2.
\]
The multiplier changes.  Thus a generic-sector affine action needs a separate
dual-label or character-invariance convention before it can be called part of
the Weyl-net symmetry theorem.
